\documentclass[11pt,a4paper]{article}
\usepackage[T1]{fontenc}
\usepackage[utf8]{inputenc}
\usepackage[margin=28mm]{geometry}
\usepackage{lmodern,amsmath,amssymb,amsthm,mathtools,microtype}
\usepackage{enumitem,xcolor,xurl,hyperref}
\hypersetup{colorlinks=true,linkcolor=blue!45!black,citecolor=blue!45!black,urlcolor=blue!45!black}
\setlength{\emergencystretch}{2em}
\newtheorem{theorem}{Theorem}[section]
\newtheorem{proposition}[theorem]{Proposition}
\newtheorem{lemma}[theorem]{Lemma}
\newtheorem{corollary}[theorem]{Corollary}
\theoremstyle{definition}\newtheorem{definition}[theorem]{Definition}
\newtheorem{example}[theorem]{Example}
\theoremstyle{remark}\newtheorem{remark}[theorem]{Remark}
\DeclareMathOperator{\Spec}{Spec}\DeclareMathOperator{\Pic}{Pic}
\DeclareMathOperator{\Gal}{Gal}\DeclareMathOperator{\Br}{Br}
\DeclareMathOperator{\Frac}{Frac}\DeclareMathOperator{\Hom}{Hom}
\DeclareMathOperator{\cd}{cd}\DeclareMathOperator{\im}{im}
\DeclareMathOperator{\Ver}{Ver}\DeclareMathOperator{\rank}{rank}
\newcommand{\F}{\mathbf F}\newcommand{\Z}{\mathbf Z}
\newcommand{\Q}{\mathbf Q}\newcommand{\OO}{\mathcal O}
\newcommand{\Ga}{\mathbf G_a}\newcommand{\Gm}{\mathbf G_m}
\newcommand{\workbenchurl}{https://github.com/KokunoYumeto/schemes-and-moduli-workbench}
\newcommand{\recordurl}{https://github.com/KokunoYumeto/mathematics-workbenches}
\title{Counterexamples to Grothendieck's cohomological-dimension conjectures\protect\\
and birational invariance of the torsion Picard scheme}
\author{}
\date{23 November 2026}
\begin{document}
\maketitle
\begin{abstract}
Grothendieck's Montpellier notes conjecture that the cohomological dimension
of an integral affine scheme is controlled by its function field, possibly
after shrinking, and state birational invariance of the torsion Picard
scheme for a regular proper variety. We construct
an excellent Dedekind domain whose function field has $2$-cohomological
dimension one while every nonempty affine open has dimension two and an
infinite-dimensional Picard group modulo two. A mixed-characteristic
regular local surface gives a second failure, detected by a wild symbol.
We recall Ishii--Yoshida's first-direct-image vanishing theorem, explain
the Picard assertion as a consequence of known finite-torsor invariance,
and give a finite-generation
criterion when the whole boundary has a quasi-projective neighborhood.
The arithmetic proof uses class field theory; the wild calculation uses
Kato's filtration in Suzuki's formulation; the Picard proof uses purity.
The mathematics and text were produced by GPT-6 Astra (ChatGPT Pro mode)
and checked by separate instances of the same model; no human has checked them.
\end{abstract}
\section{Grothendieck's questions and the results}
The distinction between a function field and an affine model is decisive
for the questions considered here. In cote 31 of the Montpellier archive,
Grothendieck proposes bounds for the cohomological dimension of an affine
scheme in terms of its function field and residue fields. On archive
sheet 84, headed \emph{Conjectures sur la dimension cohomologique des
sch\'emas affines}, he writes, for an integral affine scheme with
$\ell$ invertible,
\[\cd_\ell(X)\leq\cd_\ell(k(X)).\]
The same sheet proposes the residual-characteristic affine-open bound
answered below \cite[sheet 84]{Grothendieck31}. We use $\cd_\ell(X)$ for the least integer $d$ such
that $H^q_{\mathrm{et}}(X,\mathcal F)=0$ for every $\ell$-primary torsion
sheaf $\mathcal F$ and every $q>d$. This convention includes arbitrary
torsion sheaves, not only constant coefficients.

\begin{theorem}\label{thm:arithmetic}
There is an excellent Dedekind domain $R$, of characteristic zero with
$2\in R^\times$, such that its fraction field $K$ is algebraic over $\Q$,
\[
 \cd_2(K)=1,\qquad
 \cd_2(U)=2\quad\text{for every nonempty affine open }U\subset\Spec R.
\]
For every nonempty open $U\subset\Spec R$, the vector space
$\Pic(U)/2$ over $\F_2$ has infinite dimension. No residue field of $R$
is orderable, and
\[
 \sup_{x\in U}\{\cd_2\kappa(x)+\dim\OO_{U,x}\}=1.
\]
\end{theorem}

Thus no shrinking removes the gap between the cohomological dimensions
of this regular arithmetic curve and its generic point. Invertibility
of two, excellence and the absence of real residue fields all hold. The
mechanism is a line-bundle obstruction: Kummer theory places
$\Pic(U)/2$ inside $H^2(U,\mu_2)$.

The construction is based on the reciprocity, ray-class existence and
transfer theorems of class field theory, in the precise forms recalled
by Milne \cite{MilneCFT}. It preserves old ideal classes modulo two
while at least doubling their dimension. Simultaneously, eventual
inertness at every old prime makes the limit ring noetherian and kills
the two-primary Brauer groups of all finite extensions of its function
field. These two effects must be arranged together. Claborn proves that
every abelian group is the class group of a Dedekind domain
\cite[Theorem 7, p.~222]{Claborn}. That realization theorem does not assert
the simultaneous algebraic-number-field, residue-field and cohomological
conditions imposed here.
S.~A.~Popescu already constructs an eventually inert quadratic tower
whose integral closure of $\Z[1/2]$ is Dedekind
\cite[\S5, pp.~511--512]{Popescu2022}. His construction controls
the odd primes. The quadratic step below adds control at dyadic primes
and injectivity of old ideal classes modulo two, together with a doubling
of their dimension. These additions supply the two cohomological features.

\begin{theorem}\label{thm:wild}
Let $k=\overline{\F}_2$, let $W(k)$ be its Witt ring, and set
\[
 A=W(k)[[t]],\qquad U=\Spec A[1/t].
\]
Then $A$ is a complete excellent regular local ring of dimension two,
but $H^3(U,\F_2)\ne0$ and $\cd_2(U)=3$.
\end{theorem}

Here the omitted divisor has characteristic-zero generic point, while
the open scheme still meets characteristic two. The nonzero class is
the boundary of $\{1+2/t,t\}$ in a henselian height-one field. Kato's
ramification filtration, as computed by Suzuki
\cite[Propositions 6.1.2, 8.2.1 and 8.2.3]{Suzuki}, separates this
symbol from classes extending across the other divisor. This example
answers the residual-characteristic affine-open question in the same
notes. It does not concern the prime-to-residue-characteristic affine
vanishing theorem.

In cote 29, archive sheet 70, Grothendieck states the Picard assertion
as a theorem:
\begin{quote}
\guillemotleft\ Alors l'homomorphisme $\Pic^\tau_Y\to\Pic^\tau_X$
est un isomorphisme\guillemotright.
\end{quote}
The argument continues on sheet 71 \cite[sheets 70--71]{Grothendieck29}.
For a proper
variety $Z$, $\Pic^\tau_{Z/k}$ denotes the union of the components of
the Picard scheme whose N\'eron--Severi classes are torsion.

\begin{theorem}\label{thm:birational}
Let $f:X\to Y$ be a proper birational morphism of integral noetherian
schemes of characteristic $p>0$. If $X$ is normal and $Y$ is regular,
then $R^1f_*\OO_X=0$.

If, in addition, $X$ and $Y$ are proper varieties over an algebraically
closed field $k$, then pullback induces an isomorphism of group schemes
\[
 \Pic^\tau_{Y/k}\xrightarrow{\sim}\Pic^\tau_{X/k}.
\]
\end{theorem}

The coherent vanishing is already Ishii--Yoshida's theorem
\cite[Corollary 1.4, p.~3]{IshiiYoshida}, in the stronger form allowing
every characteristic. The Picard assertion is also a consequence of known
work: Li--Wang prove invariance of the Nori fundamental group scheme for
a normal proper source and smooth target over a perfect field
\cite[Corollary 4.6, p.~7]{LiWang}. Their statement includes the finite
commutative torsors that determine the torsion Picard scheme. We explain
this last deduction, using the published purity theorem directly.
Grothendieck's argument on sheet 71 reduces the
Picard problem to its infinitesimal structure \cite{Grothendieck29}. Points and tangent
spaces alone cannot finish that argument: the inclusion
$\alpha_p\subset\alpha_{p^2}$ is bijective on geometric points and
on tangent spaces, but its two coordinate algebras have dimensions
$p$ and $p^2$. Purity for finite flat torsors controls every finite
subgroup, including this higher infinitesimal structure. The general
purity theorem is due to Moret--Bailly; Marrama gives a published
detailed proof \cite[Theorem 6, pp.~433--434]{Marrama}.
The quoted vanishing theorem and the finite-torsor argument both
control infinitesimal information, but the Picard assertion requires
all finite subgroup schemes, not merely a tangent-space computation.

Recent work on Picard groups also studies faithfully flat changes of the
coefficient ring. The Hexagon note contributed by Mathew and produced by
ChatGPT~6 Pro proves that every invertible $R[T]$-module becomes free
after a faithfully flat extension of $R$
\cite[Theorem 1.1, p.~1]{HexagonPicard}, using Bhatt's graded results.
That theorem concerns polynomial extensions and arbitrary flat base
change; the Picard assertion here concerns a fixed proper birational map.
It is a contextual comparison and is not an input to our proofs.

Finally, for a normal proper connected variety $X$ over an algebraically
closed field, and a proper closed subset $Y$, let $Q(X,Y)$ be the
quotient of $\pi_1(X\setminus Y)$ by the closed normal subgroup
generated by wild inertia at the codimension-one components of $Y$.
This definition uses the given compactification.

\begin{theorem}\label{thm:boundary}
If $Y$ is contained in a quasi-projective open subset of $X$, then
$Q(X,Y)$ is topologically finitely generated. In particular,
$\pi_1(X\setminus\{x\})$ is topologically finitely generated when
$x$ is a closed point of codimension at least two. More generally,
if a proper connected variety $Z$ need not be normal and a closed
subset has codimension at least two on every irreducible component
and has one quasi-projective open neighborhood, the fundamental group
of its complement is finitely generated.

For normal $X$, the same conclusion for $Q(X,Y)$ holds whenever
there is a normal projective modification $X'\to X$ such that the
wild inertia at each exceptional divisorial component of the inverse
image of $Y$ has finite image in $Q(X,Y)$.
\end{theorem}

The point is the prescribed neighborhood in Raynaud's form of Chow's
lemma \cite[Chapter 4, \S4, Proposition 5, p.~30]{Raynaud}. All
exceptional divisors of the modification then lie in the open scheme,
where the covers are already \'etale. The proof therefore avoids
introducing uncontrolled boundary inertia. In the projective case,
the curve argument is already present in Drinfeld
\cite[Lemma C.2, p.~19]{Drinfeld}; Esnault--Kindler extend its
finite-field formulation to every field
\cite[Proposition 6.2, pp.~8--9]{EsnaultKindler}. Thus the projective
input is a known Lefschetz consequence. We give the deduction for an
arbitrary closed boundary, using the connectedness theorem of
Martinelli--Naranjo--Pirola \cite[Theorem 1.1, p.~31]{MNP} in the same
curve argument. The tame-curve theorem of SGA~1 supplies finitely many
generators. The neighborhood assertion then follows from the prescribed
Chow modification. The final criterion uses one finite cover that
removes the finitely many wild inertia images.

\section{Ideal classes that survive a quadratic tower}
We construct the arithmetic curve by separating a finite-group
calculation from its realization by ray class fields.

\subsection{The group calculation}
Let $V$ be a nonzero finite-dimensional $\F_2$-vector space and put
$G=(V\oplus V)\rtimes\langle\tau\rangle$, where $\tau$ exchanges the
two summands. Its elements are denoted $(x,y;\epsilon)$, with
$\epsilon\in\F_2$. Write $H=V\oplus V$ and
$\Sigma(x,y;\epsilon)=x+y$.

\begin{lemma}\label{lem:wreath}
The kernel $A_0$ of $\Sigma$ is the elementary abelian group
$\F_2\oplus V$, with $(b,u)$ corresponding to $(u,u;b)$.
The section $v\mapsto(v,0;0)$ makes its $V$-action
\[
 v\cdot(b,u)=(b,u+bv).
\]
The map of regular modules
\[
 \phi:\F_2[V]\longrightarrow A_0,\qquad e_g\longmapsto(1,g)
\]
is surjective and $V$-equivariant. It induces a surjection
$\F_2[V]\rtimes V\twoheadrightarrow G$ over $V$. Moreover,
\[
 \Ver_{G,H}(x,y;\epsilon)=(x+y,x+y).
\]
\end{lemma}
\begin{proof}
The diagonal pairs are fixed by interchange, so multiplication on
$A_0$ is addition. Conjugating $(u,u;b)$ by $(v,0;0)$ gives
$(u+bv,u+bv;b)$. The translation $e_g\mapsto e_{g+v}$ therefore
commutes with $\phi$. Its image contains $(1,0)$ and every
$(0,g)=(1,g)+(1,0)$, proving surjectivity. The images $(g,g;1)$
are conjugates of $\tau$; each has order two and its subgroup meets
$H$ only in the identity.

For transfer choose representatives $1,\tau$ for $G/H$. If
$\epsilon=0$, transfer multiplies $(x,y;0)$ by its conjugate and
gives $(x+y,x+y)$. If $\epsilon=1$, the two cosets are exchanged,
so transfer is the square, with the same value. In the alternative
coordinates $(b,u;v)\in A_0\rtimes V$, this also gives
\begin{equation}\label{eq:square}
 (b,u;v)^2=(0,bv;0).
\end{equation}
\end{proof}

For $V=\F_2$, $G$ has order eight. The two regular basis vectors
map to the two reflections $(0,0;1)$ and $(1,1;1)$; their sum is
the central diagonal element. Transfer sends either off-diagonal
element of $H$ to that central element. This is the smallest case
of the class-preserving map in the next lemma.

\subsection{One arithmetic step}
For a totally imaginary number field $F$, set
$C_F=\Pic(\OO_F[1/2])/2$. This is a finite-dimensional
$\F_2$-vector space.

\begin{lemma}\label{lem:step}
Suppose $r=\dim C_F>0$, and let $S$ be a finite set of finite
primes of $F$ containing all dyadic primes. There is a quadratic
extension $L/F$ such that every prime in $S$ is unramified and
inert, extension of ideals is injective on $C_F$, and
$\dim C_L\ge2r$.
\end{lemma}
\begin{proof}
Class field theory identifies $C_F$ with the Galois group $V$ of
the maximal elementary abelian two-extension $M/F$ unramified at
finite primes and split at dyadic primes
\cite[Chapter V, Theorems 3.5--3.6, p.~158]{MilneCFT}. Fix, over
each $\mathfrak q\in S$, one prime $\mathfrak Q$ of $M$. Prescribe
at that prime the unit squareclass giving the unramified quadratic
extension, and prescribe the squareclass of $1$ at all other
primes above $\mathfrak q$. Thus the nonsquare prescription has exactly
one chosen prime in each fiber.

The unramified quadratic class is a unit even at a dyadic prime:
an irreducible residual polynomial $T^2-T+\bar c$ lifts with
unit discriminant $1-4c$. Local squareclasses are open, by
Hensel's lemma. Weak approximation gives an element $b\in M^\times$
with all these local classes and with valuation zero at those
primes. Choose a $V$-invariant ray modulus $\mathfrak m$, supported
there and sufficiently deep to preserve the squareclasses.

Let $E/M$ be the ray class field. Since $V$ preserves the ray
subgroup, $E/F$ is Galois. In $\Gal(E/M)$ take the Artin image
of the ray class of $(b)$. Chebotarev gives a prime
$\mathfrak p\notin S$ of $F$ and a prime of $E$ whose Frobenius
is that element \cite[Chapter V, Theorem 3.23, p.~164]{MilneCFT}.
The Frobenius is trivial on $M$, so $\mathfrak p$ splits completely
in $M$. The prime $\mathfrak P$ of $M$ so obtained has ray class
$[(b)]$. Thus
\[
 \mathfrak P=(\eta),\qquad \eta/b\equiv1\pmod{\mathfrak m}
\]
for some integral $\eta\in M$. Indeed equality of ray classes
makes $\mathfrak P(b)^{-1}$ generated by an element congruent to
one, and equality of the resulting principal ideal with
$\mathfrak P$ makes its generator integral.

Adjoin all conjugate square roots:
\[
 \widetilde N=M(\sqrt{\sigma\eta}:\sigma\in V).
\]
Their squareclasses are independent: the valuation at
$\sigma\mathfrak P$ is odd only for $\sigma\eta$. Hence
$\Gal(\widetilde N/M)=\F_2[V]$. The automorphism $g\in V$
lifts by permuting the labelled generators according to
$\sqrt{\sigma\eta}\mapsto\sqrt{g\sigma\eta}$. Independence
ensures that the algebra with these generators is a field; the
permutations respect its defining relations and compose according
to the multiplication in $V$. Therefore
$\Gal(\widetilde N/F)=\F_2[V]\rtimes V$.

Quotient by $\ker\phi$ from Lemma~\ref{lem:wreath}. The resulting
field $N$ has group $G$, and its quotient $\Sigma$ is the original
extension $M/F$. Put $L=N^H$. This is quadratic over $F$.

We check all possible ramification. Outside dyadic primes and the
primes over $\mathfrak p$, the radicands are units at odd residue
characteristic, so their quadratic extensions are split or
unramified. At dyadic primes this follows from the prescription.
At $\sigma\mathfrak P$ the inertia in $\widetilde N/M$ is generated
by $e_\sigma$. Its image in $G$ is a conjugate of $\tau$ and has
trivial intersection with $H$. Since $\mathfrak p$ splits in
$M/F$, the same computation shows $N/L$ is unramified there.
Thus $N/L$ is unramified at every finite prime. There are no real
places.

For the inertness assertion, let $v\in V$ be the Frobenius at
$\mathfrak q\in S$. If $v=0$, the local Frobenius in the regular
module is one $e_g$, because exactly one conjugate radicand has
the nonsquare unramified class. Its image $(1,g)$ has order two
and maps nontrivially to $G/H$. Hence $\mathfrak q$ is inert in
$L/F$, and the prime above it splits in $N/L$. This applies in
particular at every dyadic prime.

If $v\ne0$, its decomposition group in $M/F$ is $\langle v\rangle$.
At a chosen completion exactly two conjugate radicands have the
nonsquare class, indexed by $g$ and $g+v$. The local Frobenius over
$M$ consequently maps to $\phi(e_g+e_{g+v})=(0,v)$. If $z$ is
Frobenius over $F$, then $\Sigma(z)=v$ and $z^2=(0,v)$.
Equation~\eqref{eq:square} forces the first coordinate $b$ of
$z=(b,u;v)$ to be one. It therefore maps nontrivially to $G/H$,
again proving unramified inertness.

The unramified, dyadically split elementary abelian extension
$N/L$ gives a surjection $C_L\twoheadrightarrow V\oplus V$.
Compatibility of reciprocity with extension of ideals is transfer
\cite[Chapter V, diagram and paragraph after Proposition 3.18,
p.~162]{MilneCFT}. The composite
of $C_F\to C_L$ with that surjection is therefore
$v\mapsto(v,v)$ under $C_F=V$. It is injective, proving the
injectivity of $C_F\to C_L$. The quotient of dimension $2r$
proves the required dimension bound.
\end{proof}

\subsection{The limit ring and its cohomology}
We now prove Theorem~\ref{thm:arithmetic}, including the properties
that can fail for an uncontrolled infinite algebraic tower.

\begin{proof}[Proof of Theorem~\ref{thm:arithmetic}]
Start with $F_0=\Q(\sqrt{-35})$ and
$w=(1+\sqrt{-35})/2$, so $w^2-w+9=0$ and the discriminant is
$-35$. Minkowski's bound is $(2/\pi)\sqrt{35}<4$
\cite[Theorem 4.3, pp.~70--71]{MilneANT}. The prime 2 is inert
and principal. The two primes of norm 3 are
$P=(3,w)$ and $\bar P=(3,w-1)$. Since $N(w)=9$ and $w$ is
a unit at $\bar P$, one has $(w)=P^2$. There is no element of
norm 3: for integers $x,y$,
\[
 N(x+yw)=x^2+xy+9y^2
       =(x+y/2)^2+35y^2/4.
\]
For $y\ne0$ this exceeds 3, and for $y=0$ it is a square.
Minkowski's bound now proves $\Pic(\OO_{F_0})=\Z/2$, with
generator $P$. Inverting the principal dyadic prime leaves it
unchanged, so $\dim C_{F_0}=1$.

Enumerate the rational primes with 2 first. At stage $n$ prescribe
inertness at every prime of $F_n$ above the first $n+1$ primes,
and use Lemma~\ref{lem:step} to choose $F_{n+1}$. Define
\[
 K=\bigcup_nF_n,\qquad R=\bigcup_n\OO_{F_n}[1/2].
\]
This is the integral closure of $\Z[1/2]$ in $K$. For any fixed
rational prime, splitting stops after a finite stage and all
subsequent steps are unramified. There are consequently finitely
many primes of $R$ over it, with bounded ramification indices.
For a nonzero prime $P$ of $R$, the local ring $R_P$ is the
union of the compatible number-field valuation rings. Bounded
ramification makes its value group discrete, so it is a DVR.
Each nonzero element of $R$ is contained in only finitely many
maximal ideals: its number-field norm has finitely many rational
prime divisors and each splits only finitely often.

These facts imply noetherianity directly. For an ideal $I\ne0$,
choose $a\ne0$ in $I$, and let $P_1,\ldots,P_s$ be the maximal
ideals containing $a$. At each $P_j$ choose $b_j\in I$ with
the minimum possible valuation. The ideal
$J=(a,b_1,\ldots,b_s)$ has the same localization as $I$ at every
maximal ideal, and hence $I=J$. Thus $R$ is a Dedekind domain.
This is the finite-splitting criterion of Popescu--Vraciu
\cite[pp.~40--42]{PopescuVraciu}, here proved in the form needed.
A Dedekind domain with characteristic-zero fraction field is
excellent \cite[Tag 07QS, Proposition 15.53.3(4)]{Stacks}.
The eventual-inertness construction of a Dedekind integral closure
also appears in \cite[\S5, pp.~511--512]{Popescu2022}; the
class-preserving extension lemma is the additional requirement here.

Let $T/K$ be finite. It is $EK$ for a number field $E$ containing
some $F_n$. Any element of $\Br(T)[2]$ descends, after increasing
$E$, to an element of $\Br(E)[2]$: finite presentation descends
the algebra, and a later finite stage descends the relation that
its double is zero. This class has finitely many nonzero local
invariants. Increase $n$ so that all underlying rational primes
are prescribed, replacing $E$ by $EF_n$ and the class by its
restriction if necessary. At each such prime, the local degrees in
$F_m/F_n$ are $2^{m-n}$. In a common completion, for any prime
$w_m$ of $EF_m$ over $w$ of $E$, the degree formula gives
\[
 v_2([(EF_m)_{w_m}:E_w])
 \ge m-n-v_2([E_w:(F_n)_v]).
\]
Thus, for one sufficiently large $m$, all the relevant degrees
are even, uniformly over their prolongations. Restriction multiplies
local Brauer invariants by these degrees. All archimedean places
are complex. Brauer--Hasse--Noether therefore kills the class over
$EF_m$ \cite[Chapter III, Theorem 2.1, p.~103, and Chapter VIII,
Theorem 4.2, p.~241]{MilneCFT}.
We have shown $\Br(T)[2]=0$ for every finite extension $T/K$.

Choose a maximal pro-two subgroup of the absolute Galois group of
$K$, and let $K'/K$ be its fixed field. Kummer theory and continuity
over the finite subextensions give
\[
 H^2(K',\F_2)=H^2(K',\mu_2)
 =\varinjlim_{K\subset T\subset K',\,[T:K]<\infty}\Br(T)[2]=0.
\]
The pro-prime subgroup criterion for Galois cohomological dimension
\cite[Tag 0DV9, Lemma 59.67.3]{Stacks} now gives vanishing in
every degree at least two for all two-primary torsion coefficients.
Thus $\cd_2K\le1$. There is a nonzero prime of $R$
by lying over an odd prime of $\Z[1/2]$. A uniformizer of its
DVR has odd valuation and is not a square in $K$, so Kummer
theory proves $\cd_2K=1$.

Invertible modules and their isomorphisms descend through filtered
ring colimits. Hence
\[
 \Pic(R)/2=\varinjlim_n C_{F_n}.
\]
Every transition is injective and $\dim C_{F_n}\ge2^n$.
This vector space is therefore infinite-dimensional. A nonempty
open of the noetherian curve $\Spec R$ deletes finitely many
closed points. Its Picard group is the quotient by their divisor
classes, so its reduction modulo two still has infinite dimension.
Kummer theory gives $H^2(U,\mu_2)\ne0$ for every such open.

Every closed residue field is an increasing union of finite fields
with successive degrees eventually two. Its absolute Galois group
has trivial pro-two subgroup, and its two-cohomological dimension
is zero. The generic field contains a totally imaginary number
field; thus no residue field is orderable. The generic-field bound
\[
 \cd_2(U)\le\cd_2(K)+\dim U=2
\]
follows from SGA~4, X, Corollary 4.4, p.~57 \cite{SGA4}.
Its strict-local hypothesis holds here: at a closed point the
strict local ring is a DVR with separably closed residue field
of odd characteristic, so its fraction field has two-cohomological
dimension one by X, Theorem 2.2, p.~46; for a zero-dimensional
irreducible closed subscheme the strict localization has
separably closed residue field and dimension zero. Thus the bound
applies to arbitrary two-primary torsion sheaves. Together with
the Kummer lower bound this proves $\cd_2(U)=2$. At a closed
point the residue dimension plus local-ring dimension is $0+1$;
at the generic point it is $1+0$. This proves the final supremum.
\end{proof}

\section{The torsion Picard scheme under a birational map}
We explain how finite flat torsors detect all of the torsion Picard
scheme, including its nonreduced structure. This deduces the Picard form
of the birational-invariance theorem from known finite-torsor invariance;
it makes no claim of a new birational-invariance principle.
Li--Wang's Corollary 4.6 \cite[p.~7]{LiWang} gives the latter through the
Nori fundamental group scheme. We use Moret--Bailly's purity theorem
below to make the particular Picard deduction explicit.
The coherent vanishing in
Theorem~\ref{thm:birational} is the theorem of Ishii--Yoshida
\cite[Corollary 1.4, p.~3]{IshiiYoshida}: they prove it in every
characteristic for a normal source and regular target, without an
excellence assumption. Their Theorem 1.3 treats more general locally
quasi-unmixed targets satisfying $S_3$ and pseudo-rationality in
codimension two. Applying their corollary after localization also
covers a noetherian target without a global finite-dimensional bound.

For a proper birational map $f:X\to Y$ with $Y$ normal, there is
an open $U\subset Y$, with complement of codimension at least two,
over which $f$ is an isomorphism: the inverse rational map extends
over the codimension-one points by properness. Two finite flat torsors
over a normal integral scheme that agree on a dense open agree
everywhere. Their isomorphism scheme is finite, and a rational section
of a finite scheme extends over a normal domain by integral closure
\cite[Proposition 3, p.~432]{Marrama}.

\subsection{Finite subgroups detect a proper group scheme}
For the Picard assertion we may use the full purity theorem.
Moret--Bailly's theorem says that finite locally free torsors on a
regular noetherian scheme extend uniquely across a closed subset
of codimension at least two \cite[Theorem 6]{Marrama}. It permits
noncommutative group schemes. Combining it with the uniqueness
principle proves equivalence of finite flat $G$-torsors under the
proper birational pullback $f$.

\begin{lemma}\label{lem:detect}
Let $k$ be algebraically closed of characteristic $p$, let $P$
be a proper commutative group scheme of finite type over $k$, and
let $Q\subset P$ be a closed subgroup. If every homomorphism
from a finite commutative $k$-group scheme to $P$ factors through
$Q$, then $Q=P$.
\end{lemma}
\begin{proof}
The reduced subgroup is smooth; its identity component $B$ is
an abelian variety. For $\ell\ne p$, the union of $B[\ell^n]$
is Zariski dense in $B$. Indeed its closure is a subgroup, and
a subgroup of smaller dimension has too few $\ell^n$-torsion
points: a fixed finite number of components of dimension $b$
has at most a fixed multiple of $\ell^{2nb}$ such points,
whereas $B$ of dimension $a$ has $\ell^{2na}$. Hence $B\subset Q$.

Every component of $P_{\mathrm{red}}$ contains a torsion point.
If the component of $x$ has order $n$, then $nx\in B(k)$.
Divisibility of $B(k)$ gives $b\in B(k)$ with $nb=nx$; thus
$x-b$ is torsion. Its finite cyclic subgroup is in $Q$, and
translation proves $P_{\mathrm{red}}\subset Q$.

Choose $r$ large enough to kill every nilpotent by its $p^r$-th
power. Relative Frobenius $F^r:P\to P^{(p^r)}$ then factors
through $P_{\mathrm{red}}^{(p^r)}$. Its kernel is finite and
therefore lies in $Q$. Frobenius on the smooth group
$P_{\mathrm{red}}$ is faithfully flat. For any $x\in P(T)$,
after an fppf base change there is $b\in P_{\mathrm{red}}(T)$
with $F^r(b)=F^r(x)$. Then $x-b$ is in the finite kernel.
Both terms lie in $Q$, so $P=Q$ as fppf sheaves and hence as
group schemes.
\end{proof}

\begin{proof}[Proof of the Picard assertion in Theorem~\ref{thm:birational}]
Put $P=\Pic^\tau_{X/k}$ and $Q=\Pic^\tau_{Y/k}$. These are
proper group schemes of finite type: the identity component of
the Picard scheme of a normal proper variety is proper, and the
torsion subgroup of its N\'eron--Severi group is finite
\cite[Theorem 5.4, pp.~37--38; Remarks 5.6, p.~39, and 6.19,
p.~59]{Kleiman}.
The equality $f_*\OO_X=\OO_Y$ persists after any base change
over $k$. The projection formula then shows that a line bundle
whose pullback is trivial was already trivial. Applying this
after fppf localization proves that $Q\to P$ is a monomorphism
of relative Picard sheaves. It is proper, hence a closed immersion.

For any finite commutative group scheme $H$, let $H^D$ be its
Cartier dual. Picard--torsor duality gives, naturally for $Z=X,Y$,
\[
 H^1_{\mathrm{fppf}}(Z,H^D)
   =\Hom(H,\Pic_{Z/k})
   =\Hom(H,\Pic^\tau_{Z/k})
\]
\cite[Proposition 6.2.1, pp.~50--51]{RaynaudPicard}.
The hypotheses include properness, flatness over $k$, and
$H^0(Z,\OO_Z)=k$ universally; all hold here. A $k$-point
rigidifies the relative assertion, and finite torsors over the
algebraically closed field are trivial. The final equality holds
because a finite group has torsion image in the component group.
Finite-flat torsor purity makes the two displayed Hom groups
equal under pullback. Lemma~\ref{lem:detect} now gives $Q=P$.
\end{proof}

\section{A boundary with a quasi-projective neighborhood}
We first prove the projective case and then keep the entire boundary
fixed during the passage to a projective model.
The projective input follows from the curve theorem of
Drinfeld--Esnault--Kindler \cite[Lemma C.2, p.~19]{Drinfeld}
\cite[Proposition 6.2, pp.~8--9]{EsnaultKindler}: enlarge the boundary
to a divisor containing the singular locus and its components of
codimension at least two. A cover of the original open is unramified
at each added generic divisor. The larger-boundary quotient therefore
surjects onto the original one. The proof below writes out the fixed-curve
argument in this closed-boundary formulation.

\begin{lemma}\label{lem:projective}
For a normal connected projective variety $X$ over an algebraically
closed field and a proper closed subset $Y$, the group $Q(X,Y)$ is
topologically finitely generated.
\end{lemma}
\begin{proof}
The zero-dimensional case is immediate. For a curve this is the
tame-curve theorem \cite[XIII, Corollary 2.12,
pp.~392--395]{SGA1}. Suppose $d=\dim X\ge2$, and embed
$X$ in $\mathbf P^r$. The singular locus of $X$, the components
of $Y$ of codimension at least two, and the singular and intersection
loci of its reduced divisorial components form a closed subset
$T\subset X$ of codimension at least two. Ordinary Bertini and
the dimension count for $T$ give a linear subspace $L$ of
codimension $d-1$ such that $C=X\cap L$ is smooth, avoids $T$,
and meets the boundary divisors transversally. These are finitely
many nonempty open conditions on the parameter space over the
algebraically closed field. Write $C^\circ=C\setminus Y$.

Let $V\to X\setminus Y$ be a connected finite \'etale cover
classified by $Q(X,Y)$, and let $W$ be the normalization of $X$
in its function field. It is finite over $X$, normal, integral
and projective. At a point of $C\cap Y$ the strict henselization
of $X$ is regular and the unique boundary divisor is one regular
parameter. Abhyankar's lemma
\cite[XIII, Proposition 5.2, pp.~428--429]{SGA1} trivializes the
cover after adjoining a prime-to-characteristic root of that
parameter. Each local connected component is therefore a cyclic
Kummer chart $z^e=x_1$. Transversality makes its restriction to
$C$ regular of dimension one. Over $C^\circ$ the cover is already
\'etale. Thus the base change $W\times_X C$ is a regular
projective curve.

Apply the connectedness theorem of Martinelli--Naranjo--Pirola
\cite[Theorem 1.1, p.~31]{MNP} to $W\to\mathbf P^r$.
Its image has dimension $d$, while $\dim L=r-d+1$, so the
theorem makes $W\times_X C$ connected. A connected regular
curve is integral; removing the finitely many points above $Y$
leaves a connected tame cover of $C^\circ$. The same curve
therefore tests every connected finite $Q(X,Y)$-cover. It follows
that $\pi_1^t(C^\circ)\to Q(X,Y)$ is surjective. The tame-curve
theorem makes its source finitely generated and proves the lemma.
\end{proof}

\begin{proof}[Proof of Theorem~\ref{thm:boundary}]
Let $N\subset X$ be the prescribed quasi-projective neighborhood
of $Y$. Raynaud's Chow lemma gives a blowup with center outside
$N$ whose source is quasi-projective
\cite[Chapter 4, \S4, Proposition 5, p.~30]{Raynaud}.
Properness makes the source projective. Its normalization is
finite and still projective. We obtain a normal projective
modification $q:X'\to X$ which is an isomorphism over all of
$N$. Put $Y'=q^{-1}(Y)$ and $U'=X'\setminus Y'$.
Every divisorial boundary and its valuation ring are unchanged.

The map $\pi_1(U')\to\pi_1(X\setminus Y)$ is surjective.
Indeed a connected finite \'etale cover of a normal integral
scheme is integral. Its generic fiber is a field, which stays
the same under a birational pullback; the pullback over normal
integral $U'$ is therefore connected. Tameness at the unchanged
boundary valuations gives a surjection
\[
 Q(X',Y')\twoheadrightarrow Q(X,Y).
\]
Lemma~\ref{lem:projective} proves finite generation of the
source and hence of the target. If the boundary has codimension
at least two, there are no divisorial inertia groups to quotient
out, giving the stated assertion about the full fundamental group.
One closed point has an affine neighborhood, so the point case
follows.

For a possibly nonnormal proper variety $Z$, first remove its
nilpotents, which do not change its \'etale site, and take the
disjoint union of the normalizations of its irreducible components.
This is a finite proper surjection. The inverse image of the
specified neighborhood is quasi-projective, and the inverse
image of the omitted set has codimension at least two in each
component. The normal case gives finite generation on each
component of the punctured normalization. Effective proper
descent for finite \'etale covers, in the finite-generation form
of \cite[IX, Corollary 5.2, pp.~247--248]{SGA1}, proves finite
generation downstairs. The finitely many connected components
of the descent relation contribute only finitely many extra
generators.
\end{proof}

\begin{proposition}\label{prop:finite-wild}
Let $X$ be normal proper integral and $Y\subsetneq X$ closed,
over an algebraically closed field. Suppose a normal projective
modification $X'\to X$ has finite wild inertia image in $Q(X,Y)$
at every divisorial component of $X'\setminus q^{-1}(X\setminus Y)$.
Then $Q(X,Y)$ is finitely generated. It is enough to check the
exceptional boundary divisors.
\end{proposition}
\begin{proof}
There are finitely many boundary divisors and finitely many
nonidentity elements in the union of their wild images. A
profinite group has an open normal subgroup $H$ missing all
those elements. Let $V\to X\setminus Y$ be the connected Galois
cover with quotient $Q(X,Y)/H$. Normalize $X'$ in its function
field, obtaining a normal projective variety $T$. Above
$U'=q^{-1}(X\setminus Y)$ its open $V'$ is the connected
birational pullback of $V$. As in the preceding proof,
$\pi_1(V')\to\pi_1(V)\to H$ is surjective.

For compatible prolongations of a discrete valuation under a
finite separable extension $L/K$, absolute inertia and wild
inertia satisfy
\[
 I_w=I_v\cap G_L,\qquad P_w=P_v\cap G_L.
\]
For the second identity the intersection is pro-$p$, while its
quotient in $I_w$ is a subgroup of the prime-to-$p$ tame inertia
quotient. This characterization also applies to imperfect
residue fields \cite[Tag 0BUA, Lemma 58.13.5]{Stacks}.
Every boundary divisor of $T$ lies over a boundary divisor of
$X'$, because $T\to X'$ is finite. Its wild image in $H$ is
therefore trivial, by the choice of $H$ and its normality.
The displayed surjection factors through $Q(T,T\setminus V')$.
Lemma~\ref{lem:projective} makes $H$ finitely generated, and
the finite extension $Q(X,Y)$ is then finitely generated.

The strict transform of an original boundary divisor has the
same valuation ring; its wild image was killed in the definition
of $Q(X,Y)$. Only exceptional boundary divisors need examination.
\end{proof}

\subsection{Comparison with valuation tameness}
Achinger--H\"ubner--Lara--Stix prove that $\pi_1^t(U/k)$ is
topologically finitely generated for every connected scheme $U$
of finite type over an algebraically closed field
\cite[Theorem 2.6.1, p.~23]{AHLS}. Their definition tests
valuations of residue fields over $k$, including valuations whose
centres on a proper compactification have codimension at least two
\cite[Definitions 2.2.1--2.2.2, p.~12]{AHLS}. Thus
\[
 \pi_1(U)\twoheadrightarrow Q(X,Y)
 \twoheadrightarrow\pi_1^t(U/k),\qquad U=X\setminus Y.
\]
Finite generation of the last quotient does not settle the question
for $Q(X,Y)$. The following example makes the difference explicit.

\begin{example}\label{ex:hidden-wild}
Over an algebraically closed field of characteristic $p>0$, there
is a normal projective surface $X$ with a closed point $q$ such that
\[
 Q(X,\{q\})=\pi_1(X\setminus\{q\})\simeq\Z/p\Z,
 \qquad \pi_1^t((X\setminus\{q\})/k)=1.
\]
An exceptional divisorial valuation centred at $q$ has wild inertia
image equal to $\Z/p\Z$.
\end{example}
\begin{proof}
Let $G=(\F_p,+)$ act on $S=\mathbf P^1\times\mathbf P^1$ by
$a(x,y)=(x+a,y+a)$. Every nonidentity element fixes exactly
$P=(\infty,\infty)$, so the action is free on
$S^\circ=S\setminus\{P\}$. The finite quotient $X=S/G$
is normal and projective. If $q$ is the image of $P$, then
$S^\circ\to U=X\setminus\{q\}$ is a connected finite \'etale
$G$-cover. The projective line is simply connected by
Riemann--Hurwitz, and the proper homotopy exact sequence makes
$S$ simply connected \cite[Section 58.15, Tag 0BUM]{Stacks}.
Purity identifies $\pi_1(S^\circ)$ with $\pi_1(S)$
\cite[X, Th\'eor\`eme 3.1]{SGA1}, so the Galois-cover exact
sequence gives $\pi_1(U)=G$. There is no boundary divisor on $X$,
and hence $Q(X,\{q\})=G$.

Blow up $S$ at $P$ and write $E$ for the exceptional divisor.
The action lifts. In the parameters $u=1/x$, $v=1/y$ it sends
$u$ to $u/(1+au)$ and $v$ to $v/(1+av)$. On the blowup chart
$v=ut$ the second coordinate becomes
\[
 t\longmapsto t\frac{1+au}{1+aut}.
\]
Thus $G$ fixes $E$ and acts trivially on its residue field.
For the valuation of $E$ in $k(S)/k(X)$, both decomposition
and inertia are all of $G$; since $G$ is a $p$-group,
wild inertia is also all of $G$. Its restricted valuation is divisorial on the
normal projective modification $\operatorname{Bl}_P(S)/G\to X$
and is centred at $q$. The all-valuation tame quotient must kill
this whole image, so it is trivial. This computation uses inertia,
not just a ramification index, and also allows an inseparable
residue extension.
\end{proof}

\section{A wild class on a mixed-characteristic open surface}
We prove Theorem~\ref{thm:wild} by reducing the question to the
first two graded pieces of a ramification filtration.

Write $A=W(k)[[t]]$, where $k=\overline{\F}_2$, and put
\[
 R_0=\Spec A[1/2],\quad V=\Spec A[1/(2t)],\quad
 B=k[[t]],\quad E=k((t)),\quad C=\Frac W(k).
\]
Let $D=(A_{(2)})^h$ and $L=\Frac D$. Then $D$ is an excellent
henselian DVR with residue field $E$. The coefficient sheaf
$\Lambda=\F_2$ is constant on all these schemes; on the
characteristic-zero schemes we identify it with every tensor
power of $\mu_2$.

\begin{proof}[Proof of Theorem~\ref{thm:wild}]
The regular parameters of $A$ are $2,t$. Complete noetherian
local rings are excellent. Its residue field is algebraically
closed, so the proposed residual bound equals two.
The cohomological-dimension theorems of Gabber--Orgogozo give
$\cd_2(R_0),\cd_2(V)\le2$, and $\cd_2(L)=2$
\cite[Theorems 1.2 and 6.1]{GabberOrgogozo}. More explicitly,
the equal-characteristic dimension invariant of $E=k((t))$
is one; adjoining the mixed-characteristic henselian valuation
raises it by one. Henselian invariance and the characteristic-two
field bound give $\cd_2(D)=\cd_2(E)\le1$.

The complement $z=U\setminus V$ is the point with residue field
$E$. Henselian excision and localization for $\Spec D$ give
$H^3_z(U,\Lambda)=H^2(L,\Lambda)$. Localization on $U$ thus gives
\begin{equation}\label{eq:wild-coker}
 H^2(V,\Lambda)\longrightarrow H^2(L,\Lambda)
     \longrightarrow H^3(U,\Lambda)\longrightarrow0.
\end{equation}
The same sequences with any two-primary torsion sheaf show
vanishing above degree three: the restrictions to $D,L,V$
vanish above degrees $1,2,2$, respectively. Hence $\cd_2(U)\le3$.

The point $R_0\setminus V$ is $(t)$, with residue field $C$.
Here two is invertible, so codimension-one purity applies.
Since $H^3(R_0,\Lambda)=0$, the residue sequence is exact:
\[
 H^2(R_0,\Lambda)\longrightarrow H^2(V,\Lambda)
     \longrightarrow H^1(C,\Lambda)\longrightarrow0.
\]
The class of $u\in C^\times/C^{\times2}$ is the residue of
$\{t,u\}$. Thus every class in the first image of
\eqref{eq:wild-coker} is of the form
\begin{equation}\label{eq:wild-image}
 a+\{t,u\},\qquad a\in H^2(R_0,\Lambda),\quad u\in C^\times.
\end{equation}

For $m\ge1$, let $U^mH^2(L,\Lambda)$ be generated by symbols
whose first entry belongs to $1+2^mD$, and let $U^0=H^2$.
The absolute ramification index is one and the critical index
is two. Kato's symbol calculation, in
\cite[Proposition 6.1.2(1),(3), p.~51]{Suzuki}, gives
\[
 \operatorname{gr}^0H^2(L,\Lambda)=E^\times/E^{\times2},
 \qquad \operatorname{gr}^1H^2(L,\Lambda)=\Omega_E^1.
\]
The maps are $v\mapsto\{\widetilde v,2\}$ in degree zero
and $x\,d\log y\mapsto\{1+2\widetilde x,\widetilde y\}$
in degree one. Suzuki's
Propositions 8.1.4 (pp.~64--65), 8.2.1(1),(3) (p.~65),
and 8.2.3 (p.~66) show that
$H^2(R_0,\Lambda)\to H^2(L,\Lambda)$ is injective and strict
for this filtration, with images
\[
 B^\times/B^{\times2}\subset E^\times/E^{\times2},
 \qquad \Omega_B^1\subset\Omega_E^1
\]
on these two graded pieces. The hypotheses hold because $A$
is regular, excellent and henselian, its residue field is perfect,
$-1\in A$ is the required second root of unity, and its unique
prime dividing two has regular quotient $B$.

Suppose a class \eqref{eq:wild-image} lies in $U^1$. Write
$u=2^eu_0$, with $u_0\in W(k)^\times$. Since $k$ is perfect,
divide by a square of a Teichm\"uller unit to make $u_0\in1+2W(k)$.
The degree-zero contribution of the symbol is $t^e$, whereas
that of $a$ is the class of a unit of $B$. Equality modulo
squares forces $e$ even, by the $t$-valuation. Remove that
square; the symbol is now in $U^1$, and strictness puts $a$
in $U^1$ as well. For $u=1+2b$, its degree-one contribution
is $\bar b\,d\log t$, with $\bar b\in k$. The whole degree-one
image is therefore contained in
\[
 \Omega_B^1+k\,d\log t=B\,dt+k\,t^{-1}dt.
\]
On the other hand, the class $c=\{1+2/t,t\}$ has image
$t^{-2}dt$. This does not belong to the displayed subspace:
the coefficient of $t^{-2}$ separates it. Since $k$ is perfect,
$t$ is a two-basis of both $B$ and $E$, so their absolute
differential modules are exactly $B\,dt$ and $E\,dt$; no
completed-differential convention is being used. Thus $c$ is
outside the image of $H^2(V,\Lambda)$. Its boundary in
\eqref{eq:wild-coker} is nonzero, proving the theorem.
\end{proof}

\section{Scope and remaining questions}
The arithmetic construction settles the eventual-equality assertion for
general excellent affine schemes negatively. It does not classify the
arithmetic towers for which equality does hold. In particular, growth of
finite-level class groups alone does not imply a persistent Picard class:
the injectivity of extension modulo two is essential in our construction.

The first-direct-image theorem concerns $R^1f_*\OO_X$. It does not assert
vanishing of all higher direct images for a merely normal source. The
boundary theorem requires one quasi-projective open containing the
entire boundary. For arbitrary closed boundaries on nonprojective normal
proper varieties, the argument leaves a precise question: can the images
of the exceptional wild inertia groups on a projective modification be
made finite? No such assertion is proved here. Finite generation of
$Q(X,Y)$ for an arbitrary boundary on a normal proper variety remains
unresolved here; the all-valuation tame quotient discussed in
Example~\ref{ex:hidden-wild} can be strictly smaller.

\section*{Method and records}
The mathematics and text were produced by GPT-6 Astra (ChatGPT Pro mode).
Separate instances of GPT-6 Astra, with inherited reasoning settings,
provided complete statement-only re-proofs of the arithmetic construction,
the mixed-characteristic counterexample, the four boundary assertions,
and the Picard deduction, using the stated constructions and named
published inputs. Their proofs were compared with the arguments here.
The full arguments received further targeted checks. Finite computations
were repeated using a permutation model independent of the original
group implementation. These computations do not certify the infinite
arithmetic construction. No human has checked the mathematics or text.
We dedicate this text and its new mathematical exposition to the public
domain under CC0 1.0. The programs \emph{Quadratic tower certificate}
and \emph{Wreath permutations} and their precise scopes are recorded at
\url{\workbenchurl}. Permanent records are at \url{\recordurl}.
\begin{thebibliography}{99}
\bibitem{Grothendieck31} A. Grothendieck, Montpellier manuscripts, cote 31,
notes on cohomological dimension, \url{https://grothendieck.umontpellier.fr/31.pdf}.
\bibitem{Grothendieck29} A. Grothendieck, Montpellier manuscripts, cote 29,
notes on the Picard scheme, \url{https://grothendieck.umontpellier.fr/29.pdf}.
\bibitem{MilneCFT} J. S. Milne, \emph{Class Field Theory}, version 4.03,
2020, \url{https://www.jmilne.org/math/CourseNotes/CFT.pdf}.
\bibitem{MilneANT} J. S. Milne, \emph{Algebraic Number Theory}, version 3.08,
2020, \url{https://www.jmilne.org/math/CourseNotes/ANT.pdf}.
\bibitem{PopescuVraciu} N. Popescu and C. Vraciu,
On Dedekind domains in infinite algebraic extensions,
\emph{Rend. Sem. Mat. Univ. Padova} \textbf{74} (1985), 39--44.
\bibitem{Claborn} L. Claborn, Every abelian group is a class group,
\emph{Pacific J. Math.} \textbf{18} (1966), 219--222,
\url{https://msp.org/pjm/1966/18-2/pjm-v18-n2-p04-p.pdf}.
\bibitem{Popescu2022} S. A. Popescu,
Infinite algebraic extensions of Dedekind structures,
\emph{Math. Reports} \textbf{24(74)} (2022), 505--512,
\url{https://imar.ro/journals/Mathematical_Reports/Pdfs/2022/3/14.pdf}.
\bibitem{Stacks} The Stacks Project Authors, \emph{The Stacks Project},
\url{https://stacks.math.columbia.edu}, statements cited by permanent tag.
\bibitem{SGA4} M. Artin, A. Grothendieck and J.-L. Verdier,
\emph{Th\'eorie des topos et cohomologie \'etale des sch\'emas},
SGA 4, Lecture Notes in Mathematics 269, 270, 305, Springer, 1972--1973.
\bibitem{GabberOrgogozo} O. Gabber and F. Orgogozo,
Sur la $p$-dimension des corps,
\emph{Invent. Math.} \textbf{174} (2008), 47--80.
\bibitem{Suzuki} T. Suzuki, Arithmetic duality for two-dimensional local
rings with perfect residue field,
\emph{J. Pure Appl. Algebra} \textbf{228} (2024), article 107465;
\url{https://arxiv.org/abs/2112.06491}, version 4.
\bibitem{Marrama} A. Marrama, On the restriction of finite torsors,
\emph{C. R. Math.} \textbf{364} (2026), 431--435,
\url{https://doi.org/10.5802/crmath.836}.
\bibitem{Kleiman} S. L. Kleiman, The Picard scheme, in
\emph{Fundamental Algebraic Geometry}, Mathematical Surveys and
Monographs 123, American Mathematical Society, 2005, 235--321;
\url{https://arxiv.org/abs/math/0504020}. We use the statement
numbering of this author version.
\bibitem{RaynaudPicard} M. Raynaud, Sp\'ecialisation du foncteur de Picard,
\emph{Publ. Math. IH\'ES} \textbf{38} (1970), 27--76.
\bibitem{SGA1} A. Grothendieck, \emph{Rev\^etements \'etales et groupe
fondamental}, SGA 1, Lecture Notes in Mathematics 224, Springer, 1971.
\bibitem{Raynaud} M. Raynaud, Flat modules in algebraic geometry,
\emph{Compositio Math.} \textbf{24} (1972), 11--31.
\bibitem{IshiiYoshida} S. Ishii and K.-i. Yoshida,
On vanishing of higher direct images of the structure sheaf,
\url{https://arxiv.org/abs/2410.15282}, version 3, 19 April 2025,
Corollary 1.4, p.~3.
\bibitem{LiWang} L. Li and H. Wang,
The birational invariance of fundamental group schemes,
\url{https://arxiv.org/abs/2604.23997}, version 2, 9 June 2026,
Corollary 4.6, p.~7.
\bibitem{MNP} D. Martinelli, J. C. Naranjo and G. P. Pirola,
Connectedness Bertini theorem via numerical equivalence,
\emph{Adv. Geom.} \textbf{17} (2017), 31--38,
\url{https://doi.org/10.1515/advgeom-2016-0028}.
\bibitem{Drinfeld} V. Drinfeld, On a conjecture of Deligne,
\emph{Mosc. Math. J.} \textbf{12} (2012), 515--542, 668;
\url{https://arxiv.org/abs/1007.4004}, version 7,
Lemma C.2, p.~19.
\bibitem{EsnaultKindler} H. Esnault and L. Kindler,
Lefschetz theorems for tamely ramified coverings,
\emph{Proc. Amer. Math. Soc.} \textbf{144} (2016), 5071--5080;
\url{https://page.mi.fu-berlin.de/esnault/preprints/helene/118-EK.pdf},
Proposition 6.2, pp.~8--9.
\bibitem{AHLS} P. Achinger, K. H\"ubner, M. Lara and J. Stix,
\emph{Tame fundamental groups of rigid spaces},
\url{https://arxiv.org/abs/2608.12601}, author version dated
12 August 2026, Definitions 2.2.1--2.2.2, p.~12,
and Theorem 2.6.1, p.~23.
\bibitem{HexagonPicard} \emph{Graded absolute integral closures and
Picard groups}, Hexagon preprint 2609.00105v1, 26 September 2026,
contributed by A. Mathew; provenance: ChatGPT 6 Pro,
\url{https://hexagonmath.org/2609.00105v1}.
\end{thebibliography}

\end{document}
